\chapter{תקציב ומשאבים}\label{ch:budget}

התקציב מכסה \num{18} חודשי פיתוח ופיילוט. כל הסכומים בשקלים חדשים, לפני מע״מ, ומבוססים על הנחות שמפורטות בטבלה. אלה אומדנים לצורך תכנון; הם יעודכנו בדוח ההתקדמות לפי הצעות מחיר בפועל.

\begin{table}[!htb]\centering
\caption{אומדן תקציב ל-\num{18} חודשים}\label{tab:budget}
\begin{tabularx}{\linewidth}{@{}l Y r@{}}
\toprule
\textbf{סעיף} & \textbf{הנחת חישוב} & \textbf{אלפי ₪} \\
\midrule
כוח אדם & ממוצע \num{11} משרות $\times$ \num{18} חודשים $\times$ \num{35} אלף ₪ עלות מעסיק & \num{6,930} \\
מחשוב \en{GPU} לאימון & \num{4} צמתים $\times$ \num{12} חודשים $\times$ \num{9} אלף ₪ & \num{432} \\
ענן לפיילוט & \num{6} חודשים $\times$ \num{25} אלף ₪ & \num{150} \\
רישוי תמונות אמיתיות לאימון & \num{600} אלף תמונות $\times$ \num{0.25} ₪ & \num{150} \\
ספקי חיפוש הפוך ו-\en{API} & \num{18} חודשים $\times$ \num{6} אלף ₪; רק תמונות שאינן במסלול הירוק, אומדן לבדיקה מול הצעת מחיר & \num{108} \\
ייעוץ משפטי ופרטיות & \en{GDPR}, חוק הגנת הפרטיות, סעיף \num{50} & \num{180} \\
בדיקת חדירה חיצונית & שתי בדיקות & \num{90} \\
סדנאות תקן וסבב ציבורי & \num{3} סדנאות וכנס השקה & \num{120} \\
ביטוח, משרד וציוד & \num{18} חודשים $\times$ \num{15} אלף ₪ & \num{270} \\
\midrule
\textbf{סה״כ לפני רזרבה} & & \textbf{8,430} \\
רזרבה לסיכונים (\num{12\%}) & מבוססת על חשיפת הסיכונים בפרק~\ref{ch:risk} & \num{1,012} \\
\textbf{סה״כ כולל רזרבה} & & \textbf{9,442} \\
\textbf{סה״כ מעוגל לבקשת מימון} & & \textbf{9,600} \\
\bottomrule
\end{tabularx}
\end{table}

ראו איור~\ref{fig:budget}.


\begin{figure}[!htb]
\begin{ltrdiagram}
% budget: TikZ / pgfplots source, built inside ltrdiagram (LTR)
\begin{tikzpicture}
\begin{axis}[xbar stacked, width=13.5cm, height=2.6cm, xmin=0, xmax=9600,
  ytick=\empty, axis y line=none, axis x line*=bottom,
  xlabel={\he{אלפי ₪}}, xlabel style={font=\sffamily\scriptsize},
  ticklabel style={font=\sffamily\scriptsize}, bar width=16pt,
  legend style={font=\sffamily\scriptsize,at={(0.5,1.15)},anchor=south,legend columns=4,draw=none}]
\addplot[fill=acc,draw=white] coordinates {(6930,0)};
\addplot[fill=acc!55,draw=white] coordinates {(840,0)};
\addplot[fill=acc!25,draw=white] coordinates {(660,0)};
\addplot[fill=acc2,draw=white] coordinates {(1170,0)};
\legend{\he{כוח אדם: 72\%},\he{מחשוב ונתונים: 9\%},\he{משפט, אבטחה, תקן ותפעול: 7\%},\he{רזרבה ועיגול: 12\%}}
\end{axis}
\end{tikzpicture}

\end{ltrdiagram}
\caption{התפלגות התקציב: כוח האדם הוא הרכיב הדומיננטי}\label{fig:budget}
\end{figure}

\section{בסיס האומדן ורגישות}
כל שורה בטבלה~\ref{tab:budget} נשענת על הנחה גלויה, כפי שממליץ מדריך אומדני העלות של משרד המבקר האמריקאי~\cite{gao2020}. שתי ההנחות שמשפיעות יותר מכל השאר נבדקו:
\begin{itemize}
\item \textbf{שכר:} עלייה של \num{10\%} בעלות המעסיק מוסיפה כ-\num{0.7} מיליון ₪. הרזרבה והעיגול (\num{1.17} מיליון ₪) מכסים אותה.
\item \textbf{משך:} עיכוב של שלושה חודשים בנתיב הקריטי מוסיף כ-\num{1.2} מיליון ₪ (שכר בלבד), ואז נדרשת בקשת שינוי. לכן שני הסיכונים שעלולים לגרום לכך (\req{R-01} ו-\req{R-04}) מנוהלים מהחודש הראשון.
\end{itemize}

\section{עלות יחידה: אב טיפוס מול ייצור המוני}
העלות השולית לתמונה היא המדד שקובע אם המוצר יכול לגדול. בפיילוט עלות המחשוב היא כ-\num{0.04} ₪ לתמונה. בקנה מידה של \num{50} מיליון תמונות בחודש, היעד הוא \num{0.012} ₪, בזכות שלושה מנגנונים: (\num{1}) ״מסלול ירוק״ לתמונות חתומות שמדלג על הגלאים; (\num{2}) הרצת גלאי קל ראשון, ורק תמונות שאינן חד משמעיות עוברות לגלאים הכבדים; (\num{3}) הנחת כמות אצל ספקי הענן. במנוי ממוצע של \num{3,500} ₪ בחודש ללקוח, \num{1,000} לקוחות ו-\num{50} מיליון תמונות בחודש, ההכנסה היא כ-\num{0.07} ₪ לתמונה, ולכן השוליים הגולמיים על המחשוב הם כ-\num{83\%}.

\section{מקורות מימון}
מענק חדשנות (לדוגמה, תוכנית של רשות החדשנות או \en{Horizon Europe}) בהיקף של כ-\num{40\%} (הערכה שלנו, לבדיקה מול תנאי המסלול), וסבב גיוס מוקדם (\en{Seed}) ליתרה. שתי מערכות חדשות בפיילוט יתבקשו להתחייב לרכישת מנוי לאחר הפיילוט, כהוכחת ביקוש למשקיעים.
